\chapter{Communication Complexity — Private Coin Basic}
%%% generated by the blueprint pipeline from TCSlib.CommunicationComplexity.NewmanTheorem.PrivateCoinBasic
\section{Overview}

This module introduces the private-coin protocol model for two-party communication
complexity. A private-coin protocol is a deterministic protocol in which Alice and Bob
each hold independent private randomness; correctness is measured by the probability
(over the private coins) of producing an incorrect answer on every input pair $(x, y)$.

%%%%%%%%%%%%%%%%%%%%%%%%%%%%%%%%%%%%%%%%%%%%%%%%%%%%%%%%%%%%%%%%%%%%%%%%%%%%%%%
\section{Declarations}

\begin{definition}[Private-coin protocol]
\lean{CommunicationComplexity.PrivateCoin.Protocol}
\statementsource{roughgarden-cc-bootcamp}{
  pdf-d54d402b16e2-p002-b002,
  pdf-d54d402b16e2-p011-b003
}
\statementsource{raoyehudayoff-ch03-randomized-protocols}{
  pdf-3eebfb23ab53-p003-b004,
  pdf-3eebfb23ab53-p004-b001
}
\label{CommunicationComplexity.PrivateCoin.Protocol}
\leanok
\uses{CommunicationComplexity.Deterministic.Protocol}
A \emph{private-coin protocol} with randomness spaces $\Omega_X$ and $\Omega_Y$, input
spaces $X$ and $Y$, and output type $\alpha$ is defined as a deterministic protocol
$\texttt{Deterministic.Protocol}\;(\Omega_X \times X)\;(\Omega_Y \times Y)\;\alpha$.
Alice's message function receives the pair $(\omega_x, x)$ and Bob's receives
$(\omega_y, y)$, so each player's coin flip is invisible to the other.
\end{definition}

\begin{definition}[Output node]
\lean{CommunicationComplexity.PrivateCoin.Protocol.output}
\statementsource{roughgarden-cc-bootcamp}{
  pdf-d54d402b16e2-p002-b002,
  pdf-d54d402b16e2-p004-b005
}
\label{CommunicationComplexity.PrivateCoin.Protocol.output}
\leanok
\uses{CommunicationComplexity.PrivateCoin.Protocol}
The constant \emph{output node} of a private-coin protocol that immediately returns the
value $a : \alpha$, without any further communication.
\end{definition}

\begin{definition}[Alice's message node]
\lean{CommunicationComplexity.PrivateCoin.Protocol.alice}
\statementsource{roughgarden-cc-bootcamp}{
  pdf-d54d402b16e2-p002-b002,
  pdf-d54d402b16e2-p004-b005
}
\label{CommunicationComplexity.PrivateCoin.Protocol.alice}
\leanok
\uses{CommunicationComplexity.PrivateCoin.Protocol}
Given a function $f : X \to \Omega_X \to \mathrm{Bool}$ and a continuation
$P : \mathrm{Bool} \to \mathrm{Protocol}\;\Omega_X\;\Omega_Y\;X\;Y\;\alpha$, constructs
the protocol node at which Alice sends the bit $f\,x\,\omega_x$ and execution
continues with $P(\text{bit})$.
\end{definition}

\begin{definition}[Bob's message node]
\lean{CommunicationComplexity.PrivateCoin.Protocol.bob}
\statementsource{roughgarden-cc-bootcamp}{
  pdf-d54d402b16e2-p002-b002,
  pdf-d54d402b16e2-p004-b005
}
\label{CommunicationComplexity.PrivateCoin.Protocol.bob}
\leanok
\uses{CommunicationComplexity.PrivateCoin.Protocol}
Given a function $f : Y \to \Omega_Y \to \mathrm{Bool}$ and a continuation
$P : \mathrm{Bool} \to \mathrm{Protocol}\;\Omega_X\;\Omega_Y\;X\;Y\;\alpha$, constructs
the protocol node at which Bob sends the bit $f\,y\,\omega_y$ and execution continues
with $P(\text{bit})$.
\end{definition}

\begin{definition}[Randomized execution]
\lean{CommunicationComplexity.PrivateCoin.Protocol.rrun}
\statementsource{roughgarden-cc-bootcamp}{
  pdf-d54d402b16e2-p002-b002,
  pdf-d54d402b16e2-p011-b003
}
\label{CommunicationComplexity.PrivateCoin.Protocol.rrun}
\leanok
\uses{CommunicationComplexity.Deterministic.Protocol.run, CommunicationComplexity.PrivateCoin.Protocol}
$\texttt{rrun}\;p\;x\;y\;\omega_x\;\omega_y$ executes the private-coin protocol $p$ on
inputs $x \in X$ and $y \in Y$ with Alice's private coin $\omega_x \in \Omega_X$ and
Bob's private coin $\omega_y \in \Omega_Y$, returning the output $\alpha$.  It is
defined by $p.\mathrm{run}\,(\omega_x, x)\,(\omega_y, y)$.
\end{definition}

\begin{theorem}[Randomized execution as deterministic run on paired inputs]
\lean{CommunicationComplexity.PrivateCoin.Protocol.rrun_eq}
\label{CommunicationComplexity.PrivateCoin.Protocol.rrun_eq}
\leanok
\uses{CommunicationComplexity.Deterministic.Protocol.run, CommunicationComplexity.PrivateCoin.Protocol, CommunicationComplexity.PrivateCoin.Protocol.rrun}
\difficulty{1}
Let $p$ be a private-coin protocol with randomness spaces $\Omega_X$ and $\Omega_Y$,
input spaces $X$ and $Y$, and output type $\alpha$. Then for all inputs $x \in X$ and $y
\in Y$ and all coin outcomes $\omega_x \in \Omega_X$ and $\omega_y \in \Omega_Y$, the
randomized execution of $p$ equals the deterministic execution of $p$ on the paired
inputs $(\omega_x, x)$ and $(\omega_y, y)$:
\[
\mathrm{rrun}(p, x, y, \omega_x, \omega_y) \;=\; \mathrm{run}\bigl(p, (\omega_x, x),
(\omega_y, y)\bigr).
\]
\end{theorem}

\begin{definition}[Approximate satisfaction]
\lean{CommunicationComplexity.PrivateCoin.Protocol.ApproxSatisfies}
\statementsource{roughgarden-cc-bootcamp}{pdf-d54d402b16e2-p011-b003}
\statementsource{raoyehudayoff-ch03-randomized-protocols}{pdf-3eebfb23ab53-p004-b001}
\label{CommunicationComplexity.PrivateCoin.Protocol.ApproxSatisfies}
\leanok
\uses{CommunicationComplexity.PrivateCoin.Protocol, CommunicationComplexity.PrivateCoin.Protocol.rrun}
A private-coin protocol $p$ \emph{$\varepsilon$-satisfies} a predicate
$Q : X \to Y \to \alpha \to \mathrm{Prop}$ if for every input pair $(x, y)$,
\[
  \mu\!\left\{\,\omega \in \Omega_X \times \Omega_Y \;\middle|\;
    \neg Q\,x\,y\,(p.\mathrm{rrun}\;x\;y\;\omega_1\;\omega_2)\,\right\} \;\le\; \varepsilon,
\]
where $\mu$ is the product measure on $\Omega_X \times \Omega_Y$.
\end{definition}

\begin{definition}[Approximate computation]
\lean{CommunicationComplexity.PrivateCoin.Protocol.ApproxComputes}
\statementsource{roughgarden-cc-bootcamp}{
  pdf-d54d402b16e2-p011-b003,
  pdf-d54d402b16e2-p013-b004
}
\statementsource{raoyehudayoff-ch03-randomized-protocols}{pdf-3eebfb23ab53-p004-b001}
\label{CommunicationComplexity.PrivateCoin.Protocol.ApproxComputes}
\leanok
\uses{CommunicationComplexity.PrivateCoin.Protocol, CommunicationComplexity.PrivateCoin.Protocol.rrun}
A private-coin protocol $p$ \emph{$\varepsilon$-computes} a function $f : X \to Y \to \alpha$ if
for every input pair $(x, y)$,
\[
  \mu\!\left\{\,\omega \in \Omega_X \times \Omega_Y \;\middle|\;
    p.\mathrm{rrun}\;x\;y\;\omega_1\;\omega_2 \ne f\,x\,y\,\right\} \;\le\; \varepsilon.
\]
This formalises the standard notion of bounded-error private-coin randomized
communication complexity.
\end{definition}

\begin{theorem}[Approximate computation as satisfaction of the equality predicate]
\lean{CommunicationComplexity.PrivateCoin.Protocol.ApproxComputes_eq_ApproxSatisfies}
\label{CommunicationComplexity.PrivateCoin.Protocol.ApproxComputes_eq_ApproxSatisfies}
\leanok
\uses{CommunicationComplexity.PrivateCoin.Protocol, CommunicationComplexity.PrivateCoin.Protocol.ApproxComputes, CommunicationComplexity.PrivateCoin.Protocol.ApproxSatisfies}
\difficulty{1}
Let $p$ be a private-coin protocol with randomness spaces $\Omega_X$ and $\Omega_Y$,
input spaces $X$ and $Y$, and output type $\alpha$, each randomness space carrying a
measure, let $f : X \to Y \to \alpha$ be a target function, and let $\varepsilon \in
\bbr$. Then the proposition that $p$ approximately computes $f$ within error
$\varepsilon$ coincides with the proposition that $p$ approximately satisfies, within
error $\varepsilon$, the predicate that holds of $(x, y, a)$ exactly when $a = f(x, y)$.
\end{theorem}
